\textit{Catatan edisi: halaman sumber secara eksplisit berstatus belum
selesai. Halaman itu berhenti setelah menampilkan kandidat bentuk
berderajat $n-1$ dan sebuah rantai persamaan kosong. Jika setiap
$\partial_iH_i=h$, turunan kandidat yang ditampilkan sumber justru sama
dengan $-n\omega$, bukan $\omega$; sumber juga belum membuktikan bahwa
pilihan antiturunan pada irisan-irisan koordinat menghasilkan fungsi global
yang sesuai pada himpunan terbuka sembarang. Fragmen berikut dipertahankan
sebagai fragmen solusi sumber, bukan sebagai bukti lengkap.}

Kita dapat menuliskan bentuk tersebut sebagai

\mavergleichskettedisp
{\vergleichskette
{ \omega
}
{ =} { h(x_1 , \ldots , x_n) dx_1 \wedge \ldots \wedge dx_n
}
{ } {
}
{ } {
}
{ } {
}
}
{}{}{}
dengan suatu fungsi kontinu
\maabbdisp {h} { U } {\R
} {}
dengan
\mathl{x_1 , \ldots , x_n}{} menyatakan koordinat-koordinat pada $\R^n$.
Untuk setiap tupel tetap

\mathdisp {(x_1 , \ldots , x_{i-1},x_{i+1} , \ldots , x_n)} { , }

pembatasan fungsi $h$
pada satu komponen interval dari irisan koordinat hanya bergantung pada
$x_i$; kita nyatakan pembatasan ini dengan
\mathl{h_i(x_1 , \ldots , x_{i-1},-,x_{i+1} , \ldots , x_n)}{.}
Karena $h_i$ kontinu, pada komponen interval tersebut terdapat suatu
antiturunan
\mathl{H_i(x_1 , \ldots , x_{i-1},-,x_{i+1} , \ldots , x_n)}{.}
Turunan parsialnya terhadap $x_i$ sama dengan $h_i$. Sumber kemudian
menampilkan kandidat

\mathdisp {\sum_{i = 1}^n (-1)^{i} H_i dx_1 \wedge \ldots \wedge  dx_{i-1} \wedge d x_{i+1} \wedge \ldots \wedge  dx_n} { . }

\textit{Fragmen sumber berakhir di sini. Rantai persamaan kosong pada
halaman sumber tidak dimasukkan sebagai isi matematis.}
